\chapter{Програмування}

\section{Цілі курсу}
\begin{itemize}
	\item навчитися складати алгоритми
	\item навчиться писати фрагменти коду (текст програми)
	\item навчиться читати чужий код і використовувати його.
\end{itemize}

\section{Що таке програмування?}


Слово «програма» походить від грецьких слів «pro» (перед) і «graphein» (писати). Іншими словами, це був акт написання чогось наперед або достроково.

\begin{defn}{Програмування — це}
	процес проектування, написання, тестування, вадження і підтримки комп’ютерних програм. У вужчому значенні програмування розглядається як кодування — реалізація у вигляді програми одного чи кількох взаємопов’язаних алгоритмів.
\end{defn}

У широкому сенсі процес програмування охоплює і створення, тобто розробку, алгоритмів, і аналіз потреб майбутніх користувачів програмного забезпечення. Технологія програмування постійно розвивається, що призводить до появи нових мов програмування, інструментів розробки і сфер застосування. Тому важливою ознакою програмування є регулярне опанування нових технологій.

\begin{defn}{Алгоритм — це}
	скінченна послідовність вказівок на виконання дій, спрямованих на розв'язання певної задачі.
\end{defn}
\subsection*{Основні властивості алгоритму}

\noindent \textbf{Скінченність} — алгоритм складається з обмеженої кількості кроків і має завершуватися після виконання.

\noindent \textbf{Визначеність (детермінованість)} — кожна команда алгоритму повинна бути чіткою й однозначною, без двозначності чи різних тлумачень.

\noindent \textbf{Масовість (загальність)} — алгоритм застосовується для широкого класу задач або різних вхідних даних, а не лише для одного конкретного випадку.

\noindent \textbf{Результативність} — виконання алгоритму повинно приводити до розв’язання задачі або отримання результату.

\noindent \textbf{Дискретність} — алгоритм складається з окремих кроків (операцій), які виконуються послідовно.


	\section{Структурне програмування}

	Структурне програмування --- це методологія розробки програм, яка ґрунтується 
	на використанні трьох базових конструкцій управління виконанням:
	\begin{itemize}
		\item слідування (лінійний алгоритм),
		\item розгалуження (умовний оператор),
		\item цикл (повторення).
	\end{itemize}
	
	\subsection*{Слідування (лінійний алгоритм)}
 
		\begin{defn}{Лінійні алгоритми — це}
		алгоритми, в яких дії виконуються послідовно одна за одною. 
	\end{defn}

	Це означає, що кожна команда лінійного алгоритму виконується строго в заданому порядку. Лінійні алгоритми можуть бути корисними для розв’язання простих задач, таких як знаходження суми двох чисел або переведення одиниць виміру.
	Наприклад, лінійний алгоритм для знаходження суми двох чисел може мати наступний вигляд:
	
	\textbf{Приклад:} обчислення суми двох чисел.
\tcbset{colback=blue!5!white, colframe=blue!75!black, arc=2mm, boxrule=0.8pt}

\begin{tcolorbox}[title=Алгоритм: Сума двох чисел]

		\begin{algorithmic}
			\STATE Ввести $a, b$
			\STATE Обчислити $s = a + b$
			\STATE Вивести $s$
		\end{algorithmic}

\end{tcolorbox}
	
	\subsection*{Розгалуження (умовний оператор)}
			\begin{defn}{Розгалуження — це}
	  алгоритмічна конструкція, яка дозволяє змінювати хід виконання програми залежно від істинності певної умови.
	\end{defn}
				\begin{defn}{	Умовний оператор — це}
	 конкретна реалізація розгалуження у мові програмування.
	Він перевіряє логічний вираз і, залежно від його значення (true або false), виконує одну з гілок коду.
	\end{defn}
	
  
	\textbf{Приклад:} перевірка парності числа.

	\tcbset{colback=blue!5!white, colframe=blue!75!black, arc=2mm, boxrule=0.8pt}
	
	\begin{tcolorbox}[title=Алгоритм: Перевірка парності]

			\begin{algorithmic}
				\STATE Ввести $n$
				\IF{$n \bmod 2 = 0$}
				\STATE Вивести ``Число парне''
				\ELSE
				\STATE Вивести ``Число непарне''
				\ENDIF
			\end{algorithmic}

	\end{tcolorbox}
	
	\subsection*{Цикл (повторення)}
		\begin{defn}{Цикл — це}
	алгоритмічна конструкція, яка забезпечує багаторазове виконання одних і тих самих дій (команд) доти, доки виконується певна умова або поки не буде досягнуто заданої кількості повторень..
	\end{defn}
	

	Цикли дозволяють:
		\begin{itemize}
		\item		автоматизувати повторювані обчислення;
		\item 	скоротити код, замінивши багаторазове дублювання;
		\item 	будувати алгоритми для обробки масивів даних, серій чисел, введення/виведення.
	\end{itemize}
	
 
	\textbf{Приклад:} обчислення суми чисел від $1$ до $n$.
	\tcbset{colback=blue!5!white, colframe=blue!75!black, arc=2mm, boxrule=0.8pt}

\begin{tcolorbox}[title=Алгоритм: Сума чисел від 1 до $n$]
	
		\begin{algorithmic}
			\STATE Ввести $n$
			\STATE $s := 0$
			\FOR{$i := 1$ \TO $n$}
			\STATE $s := s + i$
			\ENDFOR
			\STATE Вивести $s$
		\end{algorithmic}

\end{tcolorbox}
	
	\subsection*{	Теорема структурного програмування}
	У 1970‑х роках Едсгер Дейкстра та Ніклаус Вірт сформулювали принципи структурного програмування.
	
	Вони довели, що
	\begin{defn}{}
	
	 \textbf{будь-який алгоритм можна представити як композицію трьох базових блоків: слідування, розгалуження, цикл.} 
\end{defn}

	Логічна повнота конструкцій : слідування забезпечує лінійність і впорядкованість дій, розгалуження додає можливість приймати рішення (умовний вибір), цикл дозволяє повторювати дії довільну кількість разів. 	Разом ці три конструкції утворюють повний набір керуючих структур, достатній для опису будь-якого алгоритму.
	
	Усі сучасні мови програмування (Python, C++, Java) мають оператори для цих трьох конструкцій (if, for, while).
	
	Навіть складні алгоритми (наприклад, сортування, пошук, обробка даних) можна розкласти на комбінацію цих трьох структур.
	
	Підпрограми (функції, процедури) лише групують ці базові блоки, але не додають нових типів керування.


\section*{Контрольні питання}

\begin{enumerate}
	\item Що таке програмування? Які основні етапи воно включає?
	\item Дайте означення алгоритму. Які його основні властивості?
	\item Що таке структурне програмування? Які три базові конструкції воно використовує?
	\item Поясніть, що таке слідування (лінійний алгоритм). Наведіть приклад.
	\item Що таке розгалуження? Чим воно відрізняється від умовного оператора?
	\item Наведіть приклад умовного оператора у будь-якій мові програмування.
	\item Що таке цикл? Які основні види циклів ви знаєте?
	\item Наведіть приклад алгоритму з використанням циклу.
	\item Сформулюйте теорему структурного програмування. На чому ґрунтується її доведення?
	\item Чому саме три конструкції (слідування, розгалуження, цикл) вважаються достатніми для побудови будь-якого алгоритму?
\end{enumerate}

	
	
	
	
	

